%%%%%%%%%%%%%%%%%%%%%%%%
% IMPORTS
%%%%%%%%%%%%%%%%%%%%%%%%

\usepackage{xcolor}
\usepackage{amsmath,amssymb,amsfonts,amsthm}
\usepackage{mathtools}
\usepackage{bm}
\usepackage{dsfont}
\usepackage[]{algorithmic, algorithm}
\usepackage{enumitem}
\usepackage{natbib}
\usepackage{verbatim}


%%%%%%%%%%%%%%%%%%%%%%%%%%%%%%%%%%%%%%%%%%%%%%%%%%%%%%%%%%%
%%%%%%%%%%%%%%%%%%%%%%%%%%%%%%%%%%%%%%%%%%%%%%%%%%%%%%%%%%%

% create tables with multi-rows
\usepackage{multirow}
\usepackage{longtable}
\usepackage{booktabs}

% for links in pdf, citing urls
\usepackage[hidelinks,pagebackref]{hyperref}
% for backref in bibliography
\renewcommand*{\backref}[1]{}
\renewcommand*{\backrefalt}[4]{{\footnotesize [%
		\ifcase #1 Not cited.%
		\or Cited on page~#2%
		\else Cited on pages #2%
		\fi%
		]}}
	
\usepackage[capitalize,nameinlink]{cleveref}
\usepackage{subcaption} % for subfigures
\usepackage{caption}

\usepackage{url}
\usepackage{cite}

% use SIOPT standard; [chapter] for thesis only
\newtheorem{theorem}{Theorem}%[chapter]
\newtheorem{lemma}[theorem]{Lemma}
\newtheorem{proposition}[theorem]{Proposition}
\newtheorem{corollary}[theorem]{Corollary}
\newtheorem{definition}[theorem]{Definition}
\newtheorem{assumption}{Assumption}
\newtheorem{example}{Example}
\newtheorem{remark}{Remark}

% no indentation for new paragraph
\usepackage[parfill]{parskip}

% parskip removes space before theorem, this fixes it: (see https://tex.stackexchange.com/questions/22119/parskip-removes-spacing-before-amsthm-theorem-environments)
\makeatletter
\def\thm@space@setup{%
	\thm@preskip=\parskip \thm@postskip=0pt
}
\makeatother



%%%%%%%%%%%%%%%%%%%%%%%%
% COMMANDS
%%%%%%%%%%%%%%%%%%%%%%%%

%switch vareps and varphi
\let\temp\phi
\let\phi\varphi
\let\varphi\temp

\let\temp\varepsilon
\let\epsilon\varepsilon
\let\varepsilon\temp

% optimization and convex analysis
\DeclareMathOperator*{\argmin}{arg\,min}
\DeclareMathOperator*{\argmax}{arg\,max}
\newcommand{\prox}[1]{\mathrm{prox}_{#1}}
\newcommand{\env}[1]{\mathrm{env}_{#1}}
\DeclareMathOperator{\conv}{conv}
\newcommand{\Lmax}{L_{\max}}




% standard maths stuff
\newcommand{\iprod}[2]{\langle #1, #2 \rangle}

\newcommand{\then}{\Longrightarrow}
\newcommand{\onehalf}{\frac{1}{2}}
\newcommand{\oneover}[1]{\frac{1}{#1}}
\newcommand{\toneover}[1]{\tfrac{1}{#1}}
\newcommand{\T}{\top}

\newcommand{\R}{\mathbb{R}}
\newcommand{\Rext}{\overline{\mathbb{R}}}
\newcommand{\Rn}{\mathbb{R}^n}
\newcommand{\N}{\mathbb{N}}

\newcommand{\E}{\mathbb{E}}
\newcommand{\Prob}{\mathbb{P}}

\newcommand{\evmin}[1]{\lambda_{\min}(#1)} % min eigenvalue
\newcommand{\evmax}[1]{\lambda_{\max}(#1)} % max eigenvalue
\newcommand{\symm}{\mathbb{S}} % symm matrices
\newcommand{\SPSD}{\mathbb{S}_{+}} % symm pos semidefdef matrices
\newcommand{\SPD}{\mathbb{S}_{++}} % symm pos def matrices


\DeclareMathOperator{\Tr}{tr} % trace of matrix
\DeclareMathOperator{\dist}{dist}
\DeclareMathOperator{\dom}{dom}
\DeclareMathOperator{\diag}{diag} % extracts diagonal from matrix
\DeclareMathOperator{\Diag}{Diag} %creates matrix from vector
\DeclareMathOperator{\ri}{ri} % relative interior
\DeclareMathOperator{\intr}{int} % relative interior

\newcommand{\unitvec}[1]{\bm{\mathrm{e}}_{#1}} % unit vectors
\newcommand{\ones}[1]{\bm{1}_{#1}} % vector of ones
\newcommand{\indfun}[1]{\mathds{1}_{#1}} % indicator function 0/1
\newcommand{\elemfun}[1]{\bm{\delta}_{#1}} % indicator function 0/1
\newcommand{\Id}{\bm{\mathrm{Id}}} % identity
\newcommand{\Proj}[1]{\mathrm{Proj}_{#1}} % projection
\DeclareMathOperator{\sign}{sign} % sign

% bold math letters
\newcommand{\bw}{\bm{w}}
\newcommand{\be}{\bm{e}}
\newcommand{\bd}{\bm{d}}
\newcommand{\bg}{\bm{g}}
\newcommand{\bx}{\bm{x}}
\newcommand{\by}{\bm{y}}

% algorithms
\newcommand{\Adam}{{\texttt{Adam}}}
\newcommand{\AdamW}{{\texttt{AdamW}}}
\newcommand{\SGD}{{\texttt{SGD}}}
\newcommand{\SGDM}{{\texttt{SGD-M}}}


% project specific
\newcommand{\chinchilla}{\textbf{Chinchilla}}
\newcommand{\marin}{\textbf{Marin}}
\newcommand{\li}{\textbf{Li}}
\newcommand{\oellm}{\textbf{OpenEuroLLM}}


% \newcommand{\ttl}{{\textbf{2TL}}}
% \newcommand{\tttl}{{\textbf{3TL}}}
\newcommand{\ttl}{{\texttt{2TL}}}
\newcommand{\tttl}{{\texttt{3TL}}}

% plain text 
\newcommand{\ie}{i.e.\ }
\newcommand{\skipline}{\vspace{\baselineskip}\\}

%%%%%%%%%%%%%%%%%%%%%%%%
%%% STYLING
%%%%%%%%%%%%%%%%%%%%%%%%
\definecolor{kleinblue}{RGB}{0, 47, 167}
\definecolor{royalblue}{RGB}{0, 33, 115}
\hypersetup{
	colorlinks = True,
	linkcolor = black,
	%linkcolor = kleinblue,
	citecolor=royalblue,
	urlcolor=royalblue
} 

%%%%%%%%%%%%%%%%%%%%%%%%
%%% USEFUL
%%%%%%%%%%%%%%%%%%%%%%%%
\definecolor{todored}{RGB}{189, 30, 30}
\newcommand\revise[1]{{\color{todored}{#1}}}
\definecolor{nypink}{RGB}{216, 131, 115}
%\newcommand\todo[1]{{\color{nypink}{#1}}}
\definecolor{commblue}{RGB}{0, 66, 102}
\newcommand\comm[1]{{\color{commblue}{#1}}}

\definecolor{takeawaycolor}{RGB}{155, 157, 137}
\definecolor{caveatcolor}{RGB}{224, 122, 95}
\definecolor{StrongRed}{RGB}{230, 57, 70}
\definecolor{OtherBlue}{RGB}{29, 112, 175}

% Orange boxes for comments
\usepackage[colorinlistoftodos]{todonotes}

\newcommand{\fabian}[1]{\todo[inline,bordercolor=orange,backgroundcolor=orange!20,linecolor=orange]{\textbf{Fabian: }#1}}
\newcommand{\takeaway}[1]{\todo[inline,bordercolor=takeawaycolor,backgroundcolor=takeawaycolor!20,linecolor=takeawaycolor]{\textbf{Takeaway: }#1}}
\newcommand{\caveat}[1]{\todo[inline,bordercolor=caveatcolor,backgroundcolor=caveatcolor!20,linecolor=caveatcolor]{\textbf{Caveat: }#1}}

%%%%%%%%%%%%%%%%%%%%%%%%
